\documentclass[12pt, a4paper]{article} % Classe: article, report o book [8]

% --- Lingua e Codifica (Supporto Italiano) ---
\usepackage[utf8]{inputenc}  % Codifica input
\usepackage[T1]{fontenc}      % Codifica font
\usepackage[italian]{babel}   % Lingua italiana [2]

% --- Impaginazione e Margini ---
\usepackage[margin=2.5cm]{geometry} % Imposta margini (2.5cm su tutti i lati)

% --- Pacchetti Scientifici/Matematici ---
\usepackage{amsmath, amssymb, amsfonts, amsthm} % Matematica avanzata [2]
\usepackage{graphicx} % Per inserire immagini
\usepackage{hyperref} % Per collegamenti ipertestuali
\usepackage{booktabs} % Per tabelle professionali

% --- Formattazione Testo e Altro ---
\usepackage[document]{ragged2e}
\usepackage{xcolor}    % Per usare i colori 

% Titolo e intestazione
\title{Federico Music}
\date{20-01-2026}


\begin{document}
\maketitle

\paragraph{Nota:}
L'intervista è stata fatta in collaborazione di Giada Schiumarini
amministrativa con ruoli di analisi aziendali.

\section{Ruolo}
Federico Music è un amministrativo con esperienza di circa 2 anni all'interno
dell'azienda e si occupa di :
\begin{itemize}
      \item Esportazione degli Ordini da Access al gestionale Ad-Hoc.
      \item Controllo della correttezza dei dati esportati.
      \item Fatturazione delle commesse.
      \item Gestione delle note di credito.
      \item Gestione delle scadenze di pagamento.
\end{itemize}

\section{Situazione attuale}

\subsection{Attuale utilizzo di Access}
Federico utilizza Access per:
\begin{itemize}
      \item Esportare gli ordini creati dai commerciali verso il gestionale Ad-Hoc.
      \item Verificare i dati di pagamento e le scadenze.
      \item Controllare la correttezza dei dati inseriti nei vari campi che verranno
            esportati.
\end{itemize}

\subsection{Attuale flusso di lavoro}
L'attuale flusso di lavor di Federico parte dalla conferma d'ordine da parte
dei commerciali e si articola in:
\begin{itemize}
      \item I commerciali inseriscono le conferme d'ordine in un sistema SharePoint/Access
            (definito "Cassetto").
      \item L'amministrazione consulta il sistema due volte al giorno (8:00 e 13:30) per
            processare le nuove richieste.
      \item Controllo manuale delle conferme, delle condizioni di pagamento e dei dati
            doganali prima della validazione.
      \item Validazione in Access.
      \item Generazione automatica di un file Excel/CSV di esportazione da Access.
      \item Importazione manuale/guidata nel gestionale ERP "Ad-Hoc".
\end{itemize}

\subsection{Controlli Attuali}
Attualmente Federico effettua i seguenti controlli manuali in Access:
\begin{itemize}
      \item Controllo delle note di credito.
      \item Controllo delle scadenze di pagamento.
      \item Controllo dello stato impianti per bloccare spedizioni in caso di mancati
            pagamenti.
      \item Controlli di compliance \textit{(es. sanzioni, VIES, origine merce, ecc.)}.
\end{itemize}

\subsection{Funzionalità da mantenere}
Nell'attuale programma Access ci sono alcune funzionalità o campi dati da
mantenere:
\begin{itemize}
      \item \textbf{Campo tipologia offerta:} Distinzione tra commesse nuove e RA (Richieste
            Assistenza), RA di amplimento , RA ricambi , ecc.
      \item \textbf{Esportazione dati ad Ad-Hoc:} Funzionalità di esportazione dati verso il gestionale
            ERP Ad-Hoc.
      \item \textbf{Mantenere possibilità di modifica manuale:} Alcuni campi devono rimanere editabili per l'amministrazione come i pagamneti.
\end{itemize}

\section{Requisiti funzionali}

\subsection{Gestione ordini e commesse}
La gestione delle commesse e degli ordini deve prevedere:
\begin{itemize}
      \item \textbf{Distinzione RA vs commesse:} Necessità di distinguere chiaramente tra:
            \begin{itemize}
                  \item \textit{Commesse:} Vendita nuovi impianti.
                  \item \textit{RA (Richieste Assistenza):} Ricambi, ampliamenti, manutenzioni.
            \end{itemize}
      \item \textbf{Categorizzazione RA:} Implementazione di un campo obbligatorio per specificare la tipologia (lista di circa 10-13 voci, es: \textit{Ampliamento, Componenti, Manutenzione, Garanzia, ecc.}) per facilitare il routing verso l'Ufficio Tecnico.
      \item \textbf{Numerazione:}
            \begin{itemize}
                  \item Il numero documento, fondamentale per la gestione amministrativa , deve essere
                        univoco. Attualmente corrisponde al numero commessa per quanto riguarda le
                        commesse e al numero della RA per quanto riguarda gli ordini successivi dello
                        stesso impianto.
                  \item Possibilità di generare il numero progressivo direttamente nel nuovo
                        applicativo e passarlo ad Ad Hoc, evitando discrepanze.
            \end{itemize}
      \item \textbf{Stato "Cassetto":} Mantenimento di un'area di "staging" (conferme in attesa di processazione amministrativa). Valutare se integrarla all'interno dell'applicativo o mantenerla separata.
\end{itemize}

\subsection{Gestione pagamenti e fatturazione}
La gestione dei pagamenti e della fatturazione deve prevedere:
\begin{itemize}
      \item \textbf{Template condizioni di pagamento:} Pre-impostazione di template standard \textit{(es. 30\% ordine, 30\% merce pronta)}, ma con \textbf{totale editabilità} da parte dell'utente per gestire eccezioni.
      \item \textbf{Dati fattura:} Necessità di riportare in fattura riferimenti specifici:
            \begin{itemize}
                  \item Numero ordine cliente.
                  \item Numero ordine interno.
                  \item \textbf{Numero ordine cliente} (Dato critico spesso mancante, da rendere campo esplicito).
                  \item Riferimento DDT.
            \end{itemize}
\end{itemize}

\subsection{Alert e controlli amministrativi}
Per un supporto proattivo all'amministrazione, il sistema deve prevedere:
\begin{itemize}
      \item \textbf{Controlli geografici/sanzioni:} Sistema di \textit{alert} automatico basato sulla nazione selezionata (es. Russia, Bielorussia, Cuba)
            per segnalare embarghi o necessità di certificazioni (es. AIFEC, Dual Use).
      \item \textbf{Origine merce:} Segnalazione di potenziali conflitti di origine merce (es. prodotti \textit{Made in USA} verso paesi sanzionati).
      \item \textbf{Validazione VIES:} Verifica validità Partite IVA (specie per UE).
\end{itemize}

\section{Struttura dati e anagrafiche}

\subsection{Avvio fase amministrativa}
L'avvio della fase amministrativa deve essere subordinata alla conferma
d'ordine del cliente da parte del commerciale. Sono da gestire due casi:
\begin{itemize}
      \item \textbf{Ordine confermato dal cliente:} Il commerciale riceve esplicita conferma dal cliente,
            conferma l'ordine e l'amministrazione può procedere.
      \item \textbf{Ordine non confermato dal cliente:} In assenza di conferma scritta dal cliente il commerciale può mandare avanti l'ordine per iniziare a processarlo.
            In questo caso l'amministrazione deve essere avvisata che l'ordine non è stato confermato, ma verrà ultimato una volta confermato.
\end{itemize}

\subsection{Anagrafica clienti (livelli)}
Il sistema deve gestire una struttura a tre livelli per ogni commessa,
risolvendo l'attuale duplicazione dei clienti in Access:
\begin{enumerate}
      \item \textbf{Cliente (Buyer):} Intestatario fattura/contratto \textit{(es. Rivenditore)}.
      \item \textbf{Utilizzatore (End User):} Soggetto giuridico che usa l'impianto (fondamentale per certificati Dual Use).
      \item \textbf{Destinazione impianto:} Luogo fisico di consegna/montaggio (può differire dalla sede legale dell'utilizzatore).
\end{enumerate}
Da tenere in cosiderazione che i tre livello possano essere coincidenti
(es. Cliente che è anche utilizzatore e luogo di consegna).

\subsection{Attributi Aggiuntivi}
Alcuni attributi aggiuntivi sono fondamentali per la corretta gestione
amministrativa:
\begin{itemize}
      \item \textbf{Settore merceologico (Industry):} Campo per tracciare l'uso finale dell'impianto (es. \textit{Bakery, Pizza, Gelato}) legato alla commessa e non solo al cliente.
      \item \textbf{Lingua documenti:} La lingua dei documenti non deve dipendere rigidamente dal paese, ma essere un attributo selezionabile (es. Cliente Romeno che richiede documentazione in Italiano).
\end{itemize}

\section{Integrazioni e Codifiche}

\subsection{Integrazione con Ad-Hoc (ERP)}
L'integrazione con il gestionale Ad-Hoc deve prevedere:
\begin{itemize}
      \item \textbf{Sincronizzazione clienti:} Importazione univoca dei codici clienti da Ad-Hoc per evitare duplicati.
      \item \textbf{Codici articolo:}
            \begin{itemize}
                  \item \textit{Attuale:} Uso di codice generico "FM"("Fuori Magazzino" in quanto non esiste una corrispondenza univoca).
                  \item \textit{Desiderata:} Passaggio a codici strutturati (livello 1 distinta base) condivisi con l'Ufficio Tecnico, mantenendo però la flessibilità di modificare la descrizione in fattura per il cliente
                        (che spesso vuole descrizioni generiche tipo "Soffiante" e non codici tecnici). Faciliterebbe la costificazione , la tracciabilità e l'analisi dei dati.
            \end{itemize}
      \item \textbf{Causali contabili:}
            \begin{itemize}
                  \item \textbf{ORCAN:} Vendita con installazione (Cantiere).
                  \item \textbf{ORCVE:} Vendita pura o con installazione minima.
            \end{itemize}
\end{itemize}
\subsection{Campi fondamentali}
Per l'amministrazione sono fondamentali i seguenti campi:
\begin{itemize}
      \item Dati del Cliente \textit{(Ragione Sociale, Indirizzo, Partita IVA, \dots)}. (Da
            verificare meglio con i reparti commerciali)
      \item Numero Ordine Cliente.
      \item Numero Ordine Interno (commessa o RA). Corrisponde al numero documento in
            Ad-Hoc ed è unico per ogni ordine. \textit{(Si può rendere UNIQUE con il
                  codice)}.
      \item Partita IVA Cliente.
      \item Tipologia di vendita , se commessa nuova o RA e sua volta RA di che tipo
            \textit{(es. ampliamento, ricambi, manutenzione , ritiro , \dots)}. Da questa
            ne potrebbe dirivare il codice dell'ordine \textit{(es. COMM , VSS , RIT,
                  \dots)}.
      \item Condizioni di Pagamento , possono essere preimpostate associate a certi
            clienti. Prestare attenzione alle casistiche:
            \begin{itemize}
                  \item Pagamento a stati di avanzamento \textit{(es. 30\% anticipo, 40\% merce pronta,
                              30\% a collaudo)}.
                  \item Pagamento a 30/60/90 giorni data fattura.
                  \item Pagamento a 30/60/90 in date prestabilite \textit{(es. 30 di maggio)}.
            \end{itemize}
      \item Data Consegna Stimata.
      \item Valuta.
      \item Nazione di destinazione , con relative regole di compliance.
      \item Luogo esatto di consegna.
      \item Causali contabili \textit{ORCAN} o \textit{ORCVE}.
      \item Codice Articolo esistenti in Ad-Hoc, la descrizione deve essere editabile.
            Possibilità così di esportare con codici standard ma descrizioni
            personalizzate.
      \item Lingua Documenti , legata al singolo cliente.
      \item Tipologia di produzione del cliente (es. \textit{Bakery, Pizza, Gelato}).
\end{itemize} 

\section{Note operative per la transizione}
Federico fornisce le seguenti indicazioni operative per la transizione al nuovo
sistema:
\begin{itemize}
      \item L'amministrazione è favorevole a un periodo di utilizzo parallelo (nuovo
            software + Access) per testare i flussi.
      \item La modifica manuale dei dati \textit{(es. descrizioni in fattura, condizioni di
                  pagamento speciali)} rimane una necessità operativa inderogabile per Federico.
\end{itemize}

\end{document}